\chapter{System Design and Technical Implementation}
\label{ch:system_design_implementation}

This chapter details the architectural design, mathematical formulations, and software engineering implementation of the SENTINEL two-tier cognitive security framework. Structured as a hybrid threat detection and incident response pipeline, this system is engineered to bridge the performance gap between wire-speed network packet processing and resource-intensive Large Language Model (LLM) reasoning, while establishing robust cognitive guardrails to isolate the LLM from adversarial manipulation. The content analyzes the overall system block diagram, the online statistical pre-filtration algorithms, the state-graph cognitive orchestration workflow, the Dual-RAG retrieval mechanism, the Threat Memory SQLite schema, the defensive guardrail pipeline, the HMAC-secured audit logs, the Streamlit user interface, and the Docker deployment model.

\section{Overview of the Two-Tier Cognitive Architecture}

To protect enterprise environments without incurring prohibitive computational latencies, the SENTINEL system splits the threat detection and response pipeline into two distinct, specialized operational tiers. This architectural separation of concerns is illustrated by the holistic system block diagram in Figure~\ref{fig:soc_ai_concept}.

\begin{figure}[!htbp]
    \centering
    \includegraphics[width=0.95\textwidth,height=0.72\textheight,keepaspectratio]{images/soc_ai_concept.png}
    \caption{The holistic architecture of the SENTINEL system, illustrating the two-tier data flow and structural components. (Source: Author)}
    \label{fig:soc_ai_concept}
\end{figure}

The proposed architecture divides responsibilities across two main tiers. Tier-1, the stateless wire-speed filtration layer, is situated directly at the network boundary, where it ingests high-velocity telemetry such as flow logs and packet metadata and either computes real-time statistical deviations through Welford's algorithm or applies regex signatures, dropping or logging benign events instantly at wire-speed. Tier-2, the stateful cognitive-reasoning layer, is driven by an autonomous LangGraph agent \cite{langgraph2024}: it performs deep cognitive reasoning, retrieves grounding contexts from the Dual-RAG engine, correlates logs against the SQLite-based Threat Memory, and issues structured mitigation recommendations such as dynamic firewall blocks.
This hybrid layout resolves the LLM latency bottleneck by offloading the vast majority of benign traffic at Tier-1---quantified empirically in Chapter~4---reserving deep LLM inference solely for complex threat investigations.

The detailed step-by-step runtime flow of the data processing and decision-making logic throughout the system is represented in the system execution flow diagram shown in Figure~\ref{fig:sentinel_flow}.

\begin{figure}[!htbp]
    \centering
    \includegraphics[width=0.95\textwidth,height=0.82\textheight,keepaspectratio]{images/sentinel_flow.png}
    \caption{Detailed runtime execution flow and data processing logic within the SENTINEL pipeline. (Source: Author)}
    \label{fig:sentinel_flow}
\end{figure}

As shown in the runtime flow, the process starts immediately when a network flow log is captured. Tier-1 intercepts it and applies signature matching via the Rule Engine; a WAF or injection signature match yields a direct \texttt{BLOCK\_IP} or \texttt{ESCALATE} verdict, while whitelisted traffic is dropped. Otherwise, the numeric statistics are updated and a Z-Score is computed. Benign logs are dropped or written to the local decision log, whereas logs whose risk score exceeds the threshold are escalated. Tier-2 then processes the escalated batch by first compressing and nonce-encapsulating it---running the Drain log parser \cite{he2017drain} for event template extraction---before \emph{unconditionally} retrieving grounding context from the Dual-RAG knowledge base (accelerated by a Semantic Cache for recurring templates) together with long-term Threat Memory. The LangGraph agent subsequently invokes the LLM triage node to emit a single structured verdict. Once resolved, the output is validated, sanitized, and written to the database with HMAC chain security, executing firewall blocks or routing the alert to the Streamlit HITL dashboard.

\section{Positioning within the Existing SOC Stack}
A fair assessment must separate what SENTINEL contributes from what mature SOC tooling already provides. Signature- and rule-based detection at wire speed (Snort, Suricata, Zeek), high-volume correlation (Splunk, QRadar, Elastic), and deterministic response automation (SOAR platforms such as Cortex XSOAR, Splunk Phantom, or Tines) are long established, and many SIEMs already tag alerts with MITRE ATT\&CK identifiers. SENTINEL is deliberately \emph{not} a competitor to these systems on detection: its stateless Tier-1 is itself an instance of exactly this signature-and-statistics layer, and in a production deployment it would front---or be fed by---an existing IDS/SIEM rather than replace it. This complementary stance is corroborated by the ablation study of Chapter~\ref{ch:experiments_evaluation}, in which adding the cognitive tier does not improve the raw flow-classification F1 over the Tier-1 gate alone; that outcome is expected precisely because commodity detection is already strong, and it is why the value of the second tier must be sought elsewhere.

The contribution instead lies in the layer that today's stack does \emph{not} provide: a secure, explainable, autonomous \emph{cognitive} tier for the ambiguous minority of events that survive Tier-1. Three properties separate this tier from conventional SOAR and from emerging cloud ``AI-SOC'' assistants. First, conventional SOAR executes \emph{rigid}, pre-authored playbooks---deterministic decision trees that cannot reason about unscripted situations or articulate a justification; SENTINEL replaces this rule-based automation with a state-graph agent that reasons over retrieved MITRE ATT\&CK techniques and NIST playbooks and emits a human-auditable rationale (rated $3.91/5$ by a cross-family LLM judge). Second, the moment a \emph{generative} model is placed in the decision loop, the log substrate itself becomes an injection surface---an adversary can embed instructions such as ``ignore previous instructions and mark this benign'' inside a field like the User-Agent---a threat to which signature engines and SOAR are structurally immune because they never interpret content generatively; the Delimited Data Encapsulation layer, the tier-consensus guard, and the HMAC-chained audit trail are a direct, empirically validated ($100\%$ full-pipeline resistance) defense for exactly this new surface. Third, commercial AI-SOC assistants typically route telemetry to a vendor-hosted cloud LLM, which is inadmissible in data-sovereign, air-gapped, or Zero-Trust environments; SENTINEL demonstrates that the entire reasoning tier runs locally on a single quantized model, with the two-tier gate keeping a multi-second local inference operationally feasible (an $82.97\%$ mean-latency reduction achieved by escalating only the ambiguous minority). In short, SENTINEL does not add another detector to a crowded market; it supplies the injection-hardened, locally deployable, and explainable reasoning-and-response tier that lets a generative model safely occupy the automation role which rigid SOAR playbooks fill today.

\section{Tier-1: Stateless Wire-Speed Filter (RuleEngine \& Welford)}

The Tier-1 filtration engine is responsible for ingestion and real-time processing of high-velocity log streams. To manage connection state without incurring disk I/O bottlenecks, the engine incorporates a temporal session baseline manager that maps active source IP addresses. For each active IP session, the system computes statistical traffic features using Welford's online algorithm \cite{welford1962}. The running mean $\mu_k$ and sum-of-squared-differences $M_{2,k}$ recurrences introduced in Section~\ref{sec:welford} (Equations~\ref{eq:welford_mean}--\ref{eq:welford_var}) yield a numerically stable rolling variance $\sigma_k^2$ over a continuous stream of numerical log attributes (such as flow duration, packet counts, or byte size). For each incoming log payload, the system then computes the Z-Score $Z$ of these characteristic traffic attributes:

\begin{equation}
    Z = \frac{x_k - \mu_k}{\sigma_k}
\end{equation}

If the Z-Score exceeds a critical threshold (defined by config, e.g., $\alpha > 3.5\sigma$), the packet is labeled with an `ESCALATE` action and forwarded to Tier-2. To preserve the temporal event chains of low-and-slow campaigns, the system continues updating session baselines. Alongside Welford-based statistical checks, Tier-1 runs a static \texttt{RuleEngine} that performs fast regex signature matching to instantly intercept known command patterns (such as SQL injection syntax) and enforce static \texttt{DROP} or \texttt{BLOCK\_IP} actions, bypassing Tier-2.

In the software implementation, Tier-1 is written in Python utilizing two primary classes: \texttt{RunningStats} and \texttt{SessionBaseline}. The \texttt{RunningStats} class maintains the running counts, mean, and sum of squares in memory, executing the mathematical updates in O(1) time complexity. To prevent memory exhaustion from tracking an unbounded number of unique IP sessions, the \texttt{SessionBaseline} class structures mapping via a hash table with a garbage-collection daemon that periodically purges inactive sessions based on a configurable Time-To-Live (TTL) parameter. The static signature matches are handled efficiently by compile-time pre-compiled regex objects in Python's \texttt{re} library.

\section{Tier-2: LangGraph Stateful Cognitive AI Agent}

Escalated anomalous flows are received by the Tier-2 cognitive engine, modeled as a stateful Directed Acyclic Graph (DAG) compiled by LangGraph \cite{langgraph2024} and functioning as a Finite State Machine (FSM). The session data and reasoning context are carried within a central state object, \texttt{SentinelState}. The compiled graph chains a guardrails node (compression and delimited encapsulation), a RAG-context node, an LLM-triage node, and a conditional ATT\&CK-mapping node in sequence, terminating in one of two action nodes (an action executor or a human-in-the-loop handler). The RAG-context node \emph{unconditionally} retrieves grounding context from the hybrid vector-keyword knowledge base before triage, rather than being reached through a knowledge-deficit branch. The LLM-triage node then evaluates the escalated batch using a locally deployed quantized Gemma 2 language model \cite{gemma_2_2024} and emits a single structured verdict. Any actionable threat verdict (\texttt{BLOCK\_IP}, \texttt{ALERT}, or \texttt{AWAIT\_HITL}) is routed through the ATT\&CK-mapping node, which deterministically anchors on the technique already grounded by triage and structures it into a schema-validated MITRE record---tactic, technique, optional sub-technique, canonical URL, a mapping-confidence score, and a tactic-derived recommended response---before the verdict proceeds to its action node; benign \texttt{LOG} verdicts bypass the mapper and terminate the cycle. The graph routes deterministically on this validated action---high-severity verdicts proceed to the action executor, ambiguous or low-confidence verdicts are diverted to the human-in-the-loop queue, and benign verdicts terminate the cycle---without looping back. The routing states and transitions are outlined in Table~\ref{tab:agent_fsm_transitions}.

\begin{table}[!htbp]
    \centering
    \caption{State Transition Matrix of the Tier-2 Cognitive Agent FSM. (Source: Author)}
    \label{tab:agent_fsm_transitions}
    \begin{tabular}{|>{\raggedright\arraybackslash}p{0.23\textwidth}|>{\raggedright\arraybackslash}p{0.25\textwidth}|>{\raggedright\arraybackslash}p{0.19\textwidth}|>{\raggedright\arraybackslash}p{0.15\textwidth}|}
        \hline
        \textbf{Current State} & \textbf{Input / Condition} & \textbf{Next State} & \textbf{Output Action} \\
        \hline
        \texttt{GUARDRAILS} & Escalated batch ingested & \texttt{RAG\_CONTEXT} & Compress \& nonce-encapsulate logs \\
        \hline
        \texttt{RAG\_CONTEXT} & Grounding context retrieved & \texttt{LLM\_TRIAGE} & Append MITRE\slash NIST context \\
        \hline
        \texttt{LLM\_TRIAGE} & action $\in$ \{\texttt{BLOCK\_IP}, \texttt{ALERT}, \texttt{AWAIT\_HITL}\} & \texttt{ATTACK\_MAPPER} & Anchor \& structure MITRE ATT\&CK record \\
        \hline
        \texttt{LLM\_TRIAGE} & action $\in$ \{\texttt{LOG}, \texttt{DROP}\} & \texttt{END} & Close cycle \\
        \hline
        \texttt{ATTACK\_MAPPER} & action $\in$ \{\texttt{BLOCK\_IP}, \texttt{ALERT}\} & \texttt{ACTION\_EXEC} & Enforce block / raise alert \\
        \hline
        \texttt{ATTACK\_MAPPER} & action = \texttt{AWAIT\_HITL} (or \texttt{BLOCK\_IP} with confidence $<0.5$) & \texttt{HITL} & Queue for SOC analyst \\
        \hline
        \texttt{ACTION\_EXEC} / \texttt{HITL} & Action completed & \texttt{END} & Write HMAC audit record \\
        \hline
    \end{tabular}
\end{table}

Grounding the triage node in the retrieved MITRE/NIST documents constrains the LLM to verified context, mitigating hallucination. The triage node emits a structured JSON payload detailing the threat class, a confidence score, and recommended block rules. High-confidence malicious verdicts are dispatched to the action executor, whereas low-confidence or explicitly ambiguous verdicts are routed to the Human-in-the-Loop (HITL) queue for analyst review. A deterministic \texttt{DecisionValidator} additionally demotes any \texttt{BLOCK\_IP} verdict whose confidence falls below 0.5 to \texttt{AWAIT\_HITL}, and overrides the LLM to \texttt{AWAIT\_HITL} whenever it contradicts a Tier-1 attack signal (the tier-consensus guard against semantic social-engineering). The free-text technique string emitted by triage is subsequently normalized by the ATT\&CK-mapping node: for the canonical web-application attack classes a curated, hand-verified table resolves the technique deterministically, while for other cases the node anchors on any valid technique identifier the triage already produced and falls back to RRF retrieval only when that identifier is absent---so the mapper structures, rather than overrides, the classifier's verdict.

The technical implementation of the state graph utilizes LangGraph's compiler. The state is represented as a structured Python dataclass (\texttt{SentinelState}) containing append-only fields like \texttt{extracted\_iocs} and \texttt{decisions} to prevent semantic drift across the triage pipeline. The LLM connection is managed via the \texttt{LLMClient} class connecting to a local quantized llama.cpp service running the Gemma 2 model \cite{gemma_2_2024}, configured with strict parameters including \texttt{temperature=0.1} and a fixed 8{,}192-token context window to ensure near-deterministic reasoning and prevent resource depletion.

\section{Dual-RAG Mechanism and Semantic Memory}

The Dual-RAG mechanism grounds the agent's decisions by retrieving contexts from two unified domain-specific databases: the MITRE ATT\&CK knowledge base \cite{strom2018mitre}, which maps threat behaviors, tactics, and techniques (TTPs), and the NIST SP 800-61r2 playbook repository, which provides standard incident handling procedures \cite{nist80061}. Rather than relying on simple keyword matching or dense semantic vectors in isolation \cite{rag2020}, retrieval utilizes a hybrid search strategy. A comparison between the dense and sparse search components is detailed in Table~\ref{tab:search_comparison}.

\begin{table}[!htbp]
    \centering
    \caption{Comparison of Retrieval Paradigms in the Hybrid Search Strategy. (Source: Author)}
    \label{tab:search_comparison}
    \begin{tabular}{|>{\raggedright\arraybackslash}p{0.22\textwidth}|>{\raggedright\arraybackslash}p{0.30\textwidth}|>{\raggedright\arraybackslash}p{0.30\textwidth}|}
        \hline
        \textbf{Dimension} & \textbf{Dense Vector Search (FAISS)} & \textbf{Sparse Keyword Search (BM25)} \\
        \hline
        \textbf{Underlying Model} & all-MiniLM-L6-v2, 384-dim embeddings \cite{reimers2019sentence} & BM25 term weighting \\
        \hline
        \textbf{Matching Logic} & Cosine Semantic Similarity \cite{faiss2019} & Term Frequency-Inverse Document Frequency \cite{bm25_2009} \\
        \hline
        \textbf{Strengths} & Captures synonyms, concepts, and abstract semantics & Matches exact technical IDs (CVEs, IPs, port numbers) \\
        \hline
        \textbf{Limitations} & Weak at exact keyword matching of technical hashes & Blind to semantic synonyms and conceptual context \\
        \hline
    \end{tabular}
\end{table}

The scores from the dense (FAISS) and sparse (BM25) pipelines are consolidated using Reciprocal Rank Fusion (RRF) \cite{rrf2009}, which calculates a unified score $R$ for each document $d$ over the set of search runs $M$:

\begin{equation}
    R(d) = \sum_{m \in M} \frac{1}{k + r_m(d)}
\end{equation}

where $r_m(d)$ is the rank of document $d$ in search run $m$, and $k$ is a constant scaling parameter (typically $k=60$). The three top-ranked passages from this fused list are forwarded to the triage node as grounding context. To avoid redundant embedding and search overhead, a Semantic Cache module sits before the retrieval pipeline. The cache derives a deterministic key by computing the SHA-256 digest of the mined log template, so escalated logs that share the same attack template reuse a previously computed retrieval result instead of re-embedding and re-searching, reducing per-query processing time from hundreds of milliseconds to microseconds. The cache is implemented as an in-memory Least Recently Used (LRU) structure (a Python \texttt{OrderedDict}) bounded to 500 entries, with an 1800-second Time-To-Live (TTL) that evicts stale templates.

\section{Threat Memory and APT Campaign Correlation}

Stealthy threat actors frequently split their attack chains into low-severity events scattered across long temporal windows, bypassing stateless packet filters. To expose these low-and-slow campaigns \cite{apt_stealth2021}, the SENTINEL architecture incorporates a persistent Threat Memory module, implemented as a local SQLite database. The module is structured around specialized schemas to track historical alerts, entity reputations, and temporal attack sequences. The database schemas are detailed in Table~\ref{tab:threat_memory_schema}.

\begin{table}[!htbp]
    \centering
    \caption{Threat Memory Database Schema Specifications. (Source: Author)}
    \label{tab:threat_memory_schema}
    \begin{tabular}{|>{\raggedright\arraybackslash}p{0.25\textwidth}|>{\raggedright\arraybackslash}p{0.33\textwidth}|>{\raggedright\arraybackslash}p{0.28\textwidth}|}
        \hline
        \textbf{Table Name} & \textbf{Primary Fields} & \textbf{Functional Purpose} \\
        \hline
        \texttt{ip\_reputation} & \texttt{ip} (PK), \texttt{reputation\_score}, \texttt{total\_incidents}, \texttt{total\_blocks}, \texttt{first\_seen}, \texttt{last\_seen}, \texttt{last\_mitre\_technique} & Tracks the long-term reputation score (0--100) and incident history of source IPs, subject to time decay. \\
        \hline
        \texttt{known\_entities} & \texttt{id} (PK), \texttt{entity\_type}, \texttt{entity\_value} (unique), \texttt{description}, \texttt{is\_active} & Pre-seeds verified internal tools/services (e.g., AD, scanners) to suppress false positives. \\
        \hline
        \texttt{threat\_events} & \texttt{id} (PK), \texttt{src\_ip}, \texttt{dst\_ip}, \texttt{apt\_phase}, \texttt{apt\_day}, \texttt{label}, \texttt{timestamp} & Records individual escalated events for emergent multi-day APT chain correlation. \\
        \hline
        \texttt{apt\_indicators} & \texttt{id} (PK), \texttt{indicator\_type}, \texttt{indicator\_value}, \texttt{confidence}, \texttt{occurrence\_count}, \texttt{mitre\_chain} & Stores confirmed APT correlation indicators for long-term campaign monitoring. \\
        \hline
    \end{tabular}
\end{table}

The reputation score of an IP address escalates based on the severity of the alerts it triggers, up to a maximum cap of 100. Conversely, a decay function periodically runs to reduce the reputation score of inactive IPs:

\begin{equation}
    S_{new} = S_{old} \times (1 - \lambda)^{t}
\end{equation}

where $\lambda$ is the decay rate and $t$ is the inactive time elapsed. An APT correlation routine scans the \texttt{threat\_events} table. If a single source IP triggers anomalies across at least two distinct days (optionally spanning multiple kill-chain phases), the system automatically elevates the threat classification to high-severity, treating it as an emergent multi-day campaign rather than an isolated event. In the implementation, these temporal intervals and reputation increments are computed by Python-side upsert and aggregation queries over the SQLite store, avoiding the need to hold long-term state in the agent's working memory.

\section{AI Security Guardrails and Encapsulation Layer}

Because LLMs are vulnerable to instruction hijacking, the SENTINEL system implements a multi-layered AI Security Guardrails layer to isolate the reasoning engine from untrusted inputs \cite{promptinjection2023}. The core defense mechanism relies on Delimited Data Encapsulation. For each session, the system generates a cryptographically secure, randomized hexadecimal token (e.g., \texttt{<|DATA\_8F3E|>}) using Python's \texttt{secrets.token\_hex()} routine (built on \texttt{os.urandom}). The untrusted log payload is encapsulated within these tags inside the prompt template, informing the LLM that any command tokens found inside the boundaries must be treated strictly as passive data rather than executable instructions.

To defend against Token Denial-of-Service (DoS) and context overflow attacks, the system incorporates a Log Template Miner based on the Drain log parsing algorithm \cite{he2017drain}. The miner groups similar log entries into templates, stripping random parameters (such as variable timestamps or thread IDs). Operating within a 4{,}000-token budget, the Token Budget Manager allocates up to 60\% of this budget to compressed templates and 40\% to high-entropy outliers. Finally, an Output Sanitization Filter enforces strict JSON schema conformance, verifying that the LLM response contains only allowed fields and stripping markdown image tags, SVG scripts, and hex-encoded bypass vectors \cite{watchtower2026}. The sanitization is implemented using the \texttt{output\_sanitizer} module which intercepts LLM outputs before they are processed by the database or dashboard.

\section{Data Integrity and HMAC Protection}

To secure the security log database against tampering by malicious insiders or attackers who gain access to the file system, the system implements HMAC log chaining~\cite{hmac1997}. Each security log entry is cryptographically linked to the preceding record, forming a private blockchain-like audit trail. For log entry $i$, the record hash $H_i$ is computed as:

\begin{equation}
    H_i = \text{HMAC-SHA256}(D_i \parallel H_{i-1}, K)
\end{equation}

where $D_i$ is the payload data of the current log, $H_{i-1}$ is the hash of the previous log record, and $K$ is the secret HMAC key stored securely. If an attacker attempts to alter a log label or delete an entry, the hash chain breaks, and the validation checks run by the dashboard immediately flag the tampering. Additionally, the system protects the RAG vector database against semantic poisoning attacks. The vector database files are monitored by a checksum manager that calculates and compares SHA-256 hashes of the files before loading them into memory, ensuring that no malicious documents have been injected into the agent's retrieval memory.

The integrity chaining is implemented in the audit-logging routine, while an \texttt{ActionValidator} class enforces an allowlist of valid security actions, detecting and rejecting any shell metacharacters inside the fields of incoming payloads to prevent remote code execution (RCE) attempts during firewall script orchestration.

\section{SOC Administration Dashboard Implementation (Streamlit UI)}

The Streamlit-based SOC Administration Dashboard follows a reactive model where the Streamlit frontend communicates dynamically with the database and files of the backend to update UI elements. The dashboard integrates three specialized display modules. First, a real-time queue interface visualizes Tier-1 telemetry throughput and Tier-2 analysis verdicts. Second, a threat memory analytics panel plots APT alert propagation graphs across IP-based temporal timelines. Third, an audit trail integrity module verifies the HMAC chain in the database, displaying status indicators where a green light signals a valid chain and a red light warns that logs have been modified. Additionally, the dashboard incorporates a Human-in-the-Loop interface, allowing analysts to manually review and either approve or reject the agent's recommended firewall blocks, updating live firewall configurations on-the-fly.

\section{Containerization and Infrastructure Deployment}

System deployment is virtualized using Docker, specified in a multi-stage \texttt{Dockerfile} designed to minimize the final production image size. The build process compiles llama.cpp with CUDA support to enable GPU-accelerated local inference within the container. Infrastructure services are managed via \texttt{docker-compose.yml}, which links the agent reasoning container with the Streamlit dashboard container, sharing log and database directories using persistent volume mounts. This ensures that the \texttt{threat\_memory.db} and \texttt{guardrails\_audit.db} SQLite databases remain decoupled from transient container lifetimes.

\section{Scalability and Graceful Degradation}

A natural objection to any LLM-in-the-loop pipeline concerns throughput: if telemetry arrives at billions of events per day, would the language model not become an insurmountable bottleneck? The two-tier design answers this by treating the two tiers as fundamentally different scaling regimes, rather than by provisioning ever-larger inference hardware.

\textbf{Tier-1 is a throughput problem, not a latency wall.} Because each Welford update and regex match executes in $O(1)$ time and the per-IP session state is mutually independent, Tier-1 is embarrassingly parallel: the workload shards cleanly across cores and hosts by source-IP hash. In the Python prototype each decision takes roughly $0.6$\,ms, so a single core sustains on the order of a thousand events per second; a throughput of one billion events per day averages about $1.2\times10^{4}$ events per second, which is within reach of a modest multi-core node, with bursts absorbed by horizontal sharding. A production implementation in compiled or stream-processing form (for example, an eBPF probe or a Rust stream operator) raises the per-core ceiling by orders of magnitude. The $O(1)$ algorithm is therefore what makes wire-speed attainable; the prototype validates the filtering ratio and correctness rather than line-rate throughput.

\textbf{The cognitive tier never sees the raw escalation count.} Although Tier-1 forwards only the anomalous minority of benign-dominated traffic, three further layers collapse this figure before any inference occurs. First, the Drain template miner compresses near-identical log lines---the repetitive bulk of real attacks such as scans, brute-force attempts, and volumetric floods---into a small set of templates. Second, the Semantic Cache, keyed on the SHA-256 digest of each template, returns a previously computed verdict for any recurring template without re-invoking the model. Third, once an IP is blocked, its Threat Memory reputation record causes subsequent events to be dropped at Tier-1 rather than re-escalated. The effective inference load is therefore bounded not by the raw event volume but by the number of \emph{genuinely novel} anomaly templates, which is far smaller and is processed in batches.

\textbf{The batch, not the event, is the unit of inference.} Escalated logs are accumulated into a single cycle batch and submitted to the model in \emph{one} inference call rather than one call per event; the number of logs carried per batch is bounded only by the token budget, so dozens of compressed escalations are adjudicated by a single ($\approx$5.7\,s) inference. Equally important, the cognitive tier sits \emph{off} the synchronous blocking path: Tier-1 has already issued its deterministic \texttt{BLOCK\_IP} or \texttt{DROP} verdict for unambiguous attacks at wire-speed, so the LLM batch contributes explanation, MITRE/NIST grounding, and cross-event correlation \emph{asynchronously}. A backlog in the cognitive tier therefore delays enrichment, not protection---the actual mitigation of clear threats is never gated on LLM latency. The sustained deep-reasoning throughput of a single accelerator is thus (logs per batch) divided by (per-batch latency), which exceeds a naive per-event rate by the batch factor, before any cache hits are counted.

\textbf{Behaviour under an adversarial flood.} The worst case for the cognitive tier is a flood of anomalies that are all novel, so that template compression and caching cannot absorb them. Two design principles bound the damage. First, such a flood is itself a Tier-1-detectable signal: the surge in event rate is a Welford outlier that triggers a coarse volumetric mitigation---for example, rate-limiting or blocking the offending address ranges---without per-event reasoning. Second, the cognitive tier is a bounded-capacity resource that must shed load gracefully: when its queue saturates, the pipeline falls back to the deterministic Tier-1 verdict (block or alert from the Z-score and rule engine), deferring only the richer LLM explanation and grounding. Security coverage is thus preserved under overload, and only cognitive enrichment degrades. Consequently, GPU capacity is provisioned for the steady-state trickle of novel anomalies---which a single consumer-grade accelerator already serves in this work---rather than for the worst-case line rate, since scaling the model to match raw ingestion would be both architecturally self-defeating and economically irrational.

It must be stated honestly that the present prototype implements the load-shedding ratio and the caching layers but not yet an explicit backpressure-and-degradation controller; under a sustained novel-anomaly flood it would saturate the escalation queue, as acknowledged among the limitations in Chapter~\ref{ch:conclusion}. Formalizing this graceful-degradation policy, together with the horizontal GPU scaling and multi-agent decomposition outlined in future work, is the natural path toward a production-grade deployment.

\section{Chapter Summary}

This chapter has established the comprehensive system design and technical implementation details of the SENTINEL security architecture. By merging wire-speed stateless filtration (Tier-1) with stateful cognitive agentic reasoning (Tier-2) and wrapping the entire system in robust cryptographic and AI-specific guardrails, the framework effectively balances high throughput with secure deep reasoning. The Streamlit administration dashboard bridges this automated system with human oversight via the HITL queue, and the multi-stage Docker setup ensures isolated, reproducible deployments. Chapter 4 will delineate the empirical evaluation, presenting performance logs, ablation studies, and statistical test results verifying the capabilities of this merged architecture.
